\section{Заключение}
\label{sec:Chapter5} \index{Chapter5}

\subsection{Решённая задача}

В работе исследована проблема качества Bottom-Level Acceleration
Structure, построенной вендорским DXR/Vulkan-билдером по входным
данным современного игрового движка. На основе этого сформулирована и
экспериментально проверена гипотеза о том, что выделение отдельного BLAS под каждую материальную подгруппу, разбросанную по всему объёму модели, порождает огромные BLAS-инстансы,
которые создают сильные перекрытия в структуре верхнего уровня,
что может существенно ухудшать время трассировки лучей. Эта гипотеза, выдвинутая на основе
эмпирических наблюдений на сценах
DagorEngine, получила в данной работе
\emph{количественное} подтверждение: на реалистичном разбиении модели
по материалам реализация предложенного метода даёт снижение времени
трассировки теней в $2{,}7$ раза, а на синтетическом худшем
разбиении --- до $8{,}6$ раза.

\subsection{Полученные результаты}

\begin{enumerate}
    \item Описана полная цепочка преобразований
      от ресурса модели до массива \texttt{RaytraceGeometryDescription[]}
      и итогового BLAS в DagorEngine. Идентифицирована доступная точка
      разработческого вмешательства, совместимая с контрактом DXR --
      стадия формирования \texttt{RElem}-массива перед вызовом
      \texttt{bvh::add\_mesh()}.
    \item Разработан алгоритм глобальной пространственной кластеризации:
      объединение треугольников всех \texttt{RElem}-ов модели,
      глобальная Morton-сортировка по центроидам, разбиение на
      пространственно-плотные кластеры с автоматической
      корректировкой их числа под аппаратный лимит 16-битного
      индекса. Линейно-логарифмическая сложность $O(N\log N)$,
      пиковая память $O(N)$, корректность доказана по построению.
    \item Реализация выполнена в семпле \texttt{testGI} движка
      DagorEngine. Код алгоритма расположен в
      \texttt{test\_app.cpp}, а управление вынесено в опцию
      \texttt{bvhMergeAllRelems}. Метод встраивается как hook
      перед штатными вызовами BVH-подсистемы и не требует
      модификации ассет-пайплайна.
    \item На целевой сцене из 16 инстансов модели метод даёт до $8{,}6$-кратного ускорения основного
      прохода GPU-трассировки теней на синтетическом случайном split:
      \texttt{rtsm::do\_trace} снижается с $9{,}15$~мс до $1{,}07$~мс; на реалистичном материальном split ---
      \aptimes{2{,}7} (с 2,75 до 1,01~мс).
    \item Визуальное подтверждение --- захваты
      Acceleration-Structure в Nsight~Graphics. В режиме раскраски
      по инстансам вместо «шумовой» картины перекрывающихся
      BLAS-инстансов появляются крупные пространственно-связные
      цветные области, соответствующие отдельным BLAS-кластерам.
\end{enumerate}

\subsection{Соответствие критериям приёмки}

В постановке задачи сформулирован
набор формальных критериев; ниже приведено их соответствие выполненной работе:

\begin{enumerate}
    \item Реализован метод как опциональный режим инициализации
      BVH-контекста, включаемый одной строкой в
      \texttt{settings.blk} и собирающийся штатным
      jam-скриптом DagorEngine. \textbf{Выполнено.}
    \item На целевой сцене зафиксировано ускорение
      до \aptimes{8{,}6} на худшем случае и \aptimes{2{,}7} на реалистичном случае против критерия не менее \aptimes{2}.
      \textbf{Выполнено.}
    \item Визуальная картина Acceleration~Structure
      демонстрирует пространственно-связные области,
      подтверждая работу алгоритма. \textbf{Выполнено.}
    \item Тени до и после применения метода визуально
      неотличимы; множество треугольников сохраняется по
      построению. \textbf{Выполнено по построению.}
\end{enumerate}

\subsection{Научная новизна и практическая значимость}

\textbf{Научная новизна} работы заключается в следующих
результатах:
\begin{itemize}
    \item Формализована проблема деградации TLAS: разбросанные по объёму модели элементы одного материала порождают огромные BLAS, которые при сборке TLAS образуют сильно перекрывающиеся листовые узлы, приводя к огромному числу ложных пересечений и резкому падению производительности.
    \item Предложен метод глобальной пространственной кластеризации ассетов,
      занимающий уникальную нишу:
      \emph{вход билдера + DXR-совместимость + сохранение
      визуальной эквивалентности} (см.
      сводную таблицу сравнения подходов). Существующие подходы либо
      работают внутри билдера (HLBVH, SBVH), либо разрезают
      отдельные треугольники без учёта поэлементной структуры
      ассета (Ganestam-Doggett).
    \item Эффективность метода продемонстрирована
      \emph{не на симуляторе и не на офлайн-метрике}, а через
      реальный замер GPU-времени в работающем семпле
      production-движка. Восьмикратное снижение
      \texttt{rtsm::do\_trace} --- редкий по величине эффект для
      одной локальной оптимизации, что подтверждает существенное
      влияние проблемы пространственной локализации в DXR-BVH.
\end{itemize}

\textbf{Практическая значимость}: метод применяется как
runtime-опция инициализации BVH-контекста для динамических
ассетов и интегрируется в существующий код движка без изменения
формата хранения мешей. Это делает его пригодным для применения
в production-проектах, использующих большие массивы готовых
ассетов, в которых перезапекание невозможно или слишком дорого.

\subsection{Направления дальнейшего развития}

\begin{enumerate}
    \item \textbf{Сохранение материальной атрибуции.} Текущая
      реализация теряет материальный идентификатор отдельных
      треугольников, что приемлемо для трассировки теней RTSM,
      но не для подсистем, использующих hit-shader (отражения). Решение: для каждого кластера хранить
      side-buffer с соответствием «локальный треугольник --- материал»
      и обращаться к нему из hit-shader-а через
      индекс вклада инстанса в DXR.
    \item \textbf{Автоматический подбор $K$.} Текущее значение
      $K$ задаётся вручную в~\texttt{settings.blk}.
      Эвристика подбора по характеристикам модели (число
      треугольников, диагональ AABB, отношение AABB к
      пространству сцены) позволит избежать ручной настройки
      под каждую модель.
    \item \textbf{Альтернативные стратегии кластеризации.}
      Глобальная Morton-кластеризация --- простой и хорошо
      работающий baseline. На моделях с сильно неоднородным
      распределением треугольников более эффективными могут
      оказаться k-means по центроидам или octree-кластеризация
      с учётом SAH-cost внутри кластера.
    \item \textbf{Полная интеграция в ассет-пайплайн.} Текущая
      реализация выполняет split каждый раз при инициализации
      BVH-контекста; для production-проектов целесообразно
      перенести Morton-сортировку и разбиение в офлайн-стадию
      запекания ассета.
    \item \textbf{Валидация на дополнительных сценах.} Текущий
      результат получен на одной целевой сцене с одной моделью.
      Расширение замеров на production-сцены DagorEngine
      позволит оценить применимость метода в
      реальных условиях и на более сложных наборах данных.
\end{enumerate}

\subsection{Итог}

Поставленная задача --- разработать и экспериментально оценить
метод предобработки меша, повышающий качество BLAS, построенного
вендорским билдером, без модификации драйвера и с измеримым
ускорением реальной GPU-трассировки --- решена. В результате разработанный
алгоритм глобальной пространственной кластеризации, его реализация в
семпле \texttt{testGI} DagorEngine и проведённые замеры образуют
законченный набор результатов для дальнейшего исследования и
переноса метода в production-окружение.

Все исходные коды реализации, тестовая модель, скриншоты
Nsight-захватов и сопровождающие материалы опубликованы в
репозитории работы~\cite{thesisRepo}.
